% Rozdział 5 — Analiza wyników i ocena modeli
% Liczby z kanonicznych artefaktów reports/ (FINAL_REPORT.md), stan 2026-07-07.
% Nowe cytowania: \cite{hardt2016}, \cite{bird2020}, \cite{grinsztajn2022}.

\section{Metryki oceny jakości modeli klasyfikacyjnych}

\subsection{Dokładność, precyzja, czułość, F1 --- i dlaczego nie są metrykami wiodącymi}

Przy rozkładzie klas 78/22 dokładność jest myląca: klasyfikator zwracający
zawsze klasę większościową osiąga 78\% przy zerowej użyteczności biznesowej.
Precyzja, czułość i~F1 zależą od progu klasyfikacyjnego, który w~niniejszym
systemie nie jest stały (progi kosztowe per model, sekcja~4.7) --- ich wartości
opisują więc punkt pracy, nie model. Z~tych powodów metryki progowe pełnią
w~pracy rolę pomocniczą (czułość i~odsetek fałszywych alarmów raportowane są
w~analizie reguł decyzyjnych, sekcja~5.4, oraz w~audycie sprawiedliwości,
sekcja~5.5), a~porównanie modeli opiera się na metrykach niezależnych od progu.

\subsection{Krzywa ROC i pole pod krzywą (AUC-ROC)}

AUC --- prawdopodobieństwo, że losowy defaultujący otrzyma wyższe PD niż losowy
spłacający --- jest niezależne od progu i~proporcji klas, co czyni je standardem
porównawczym w~credit scoringu [10]. Raportowane są także pochodne: Gini
($2\cdot\mathrm{AUC}-1$) oraz statystyka Kołmogorowa--Smirnowa (maksymalna
separacja dystrybuant wyników obu klas), tradycyjnie używane w~praktyce
bankowej.

\subsection{Wynik Briera i macierz pomyłek}

Ponieważ trajektoria PD wymaga skali absolutnej (rozdz.~4.6), obok zdolności
rankingowej oceniana jest jakość kalibracji --- wynikiem Briera (średni kwadrat
odchylenia prognozy od wyniku binarnego; niższy = lepszy). Macierze pomyłek
prezentowane są przy progach kosztowych z~sekcji~4.7, czyli w~faktycznym punkcie
pracy systemu.

\section{Wyniki modeli --- zestawienie porównawcze}

Wszystkie wyniki dotyczą zamrożonego zbioru testowego (6~000 klientów, w~tym
1~327 defaultujących) i~modeli po kalibracji izotonicznej.

\begin{table}[htbp]
\centering
\caption{Metryki pięciu modeli W3 (zbiór testowy).}
\label{tab:metryki}
\begin{tabular}{lrrrrr}
\toprule
Model & AUC & Gini & KS & Brier & Próg kosztowy \\
\midrule
Random Forest & 0,7741 & 0,548 & 0,408 & 0,1374 & 0,145 \\
XGBoost       & 0,7761 & 0,552 & 0,423 & 0,1360 & 0,165 \\
LightGBM      & 0,7767 & 0,553 & 0,417 & 0,1363 & 0,160 \\
\textbf{CatBoost} & \textbf{0,7793} & \textbf{0,559} & 0,415 & \textbf{0,1357} & 0,160 \\
LSTM          & 0,7614 & 0,523 & 0,399 & 0,1388 & 0,155 \\
\bottomrule
\end{tabular}
\end{table}

\textbf{Random Forest} (5.2.1) osiąga AUC 0,7741 --- najniższe wśród modeli
drzewiastych, przy najlepszej stabilności rankingu cech; jego ważności
(dominacja \texttt{PAY\_max}, \texttt{late\_count},
\texttt{recent\_pay\_status}) są punktem odniesienia dla analizy
interpretowalności. \textbf{XGBoost} (5.2.2, AUC 0,7761) i~\textbf{LightGBM}
(5.2.3, AUC 0,7767) osiągają wyniki praktycznie nierozróżnialne --- zgodnie
z~oczekiwaniem dla dwóch implementacji tej samej idei wzmacniania gradientowego
na wspólnej przestrzeni cech. \textbf{CatBoost} (5.2.4) jest najlepszy w~obu
wymiarach --- dyskryminacji (AUC 0,7793) i~kalibracji (Brier 0,1357).
\textbf{LSTM} (5.2.5, AUC 0,7614) zamyka stawkę: na tabelarycznych danych
kredytowych o~zaledwie trzech krokach czasowych surowa sekwencja nie
rekompensuje braku ręcznie zaprojektowanych cech agregujących --- wynik spójny
z~literaturą benchmarkową wskazującą przewagę metod drzewiastych na danych
tabelarycznych [10], [11], \cite{grinsztajn2022}. Model sekwencyjny wnosi
jednak dwie unikalne własności, omówione w~sekcjach 5.4 (jedyny model, dla
którego reguła monitorująca wygrywa również na czułości) i~5.5 (najbliższy
parytetu demograficznego --- jako jedyny nie otrzymuje na wejściu żadnych cech
demograficznych).

\begin{figure}[htbp]
\centering
\includegraphics[width=.8\textwidth]{roc_comparison_w3}
\caption{Krzywe ROC pięciu modeli W3 na zbiorze testowym. Źródło: opracowanie własne.}
\label{fig:roc}
\end{figure}

\section{Wzajemne porównanie pięciu modeli}

Różnice AUC między modelami drzewiastymi (0,0026--0,0052) są małe, dlatego
zbadano ich istotność bootstrapem: 40~repróbkowań zbioru testowego ze
zwracaniem, w~wariancie sparowanym dla różnic par modeli.

\begin{table}[htbp]
\centering
\caption{Bootstrap AUC (40 powtórzeń; 95\% przedziały percentylowe).}
\label{tab:bootstrap}
\begin{tabular}{lrl}
\toprule
Model & AUC & 95\% CI \\
\midrule
Random Forest & 0,7741 & $[0{,}765;\,0{,}786]$ \\
XGBoost       & 0,7761 & $[0{,}767;\,0{,}789]$ \\
LightGBM      & 0,7767 & $[0{,}766;\,0{,}790]$ \\
CatBoost      & 0,7793 & $[0{,}770;\,0{,}793]$ \\
LSTM          & 0,7614 & $[0{,}753;\,0{,}772]$ \\
\bottomrule
\end{tabular}
\end{table}

Sparowane przedziały różnic prowadzą do trzech wniosków: (1)~przewaga CatBoost
nad lasem losowym ($+0{,}0051$) i~XGBoost ($+0{,}0035$) jest rozróżnialna
(przedziały nie zawierają zera), nad LightGBM --- nie; (2)~różnice w~trójce
RF/XGBoost/LightGBM mieszczą się w~szumie próbkowania i~należy je raportować
jako porównywalne; (3)~przewaga wszystkich modeli drzewiastych nad LSTM
($+0{,}013$ do $+0{,}018$) jest wyraźna i~stabilna. Bootstrap repróbkuje
wyłącznie zbiór testowy przy ustalonych modelach, nie obejmuje więc wariancji
treningu.

\section{Reguła statyczna (W3) vs reguła monitorująca (W0..W3) --- weryfikacja empiryczna Wariantu B}

Właściwym eksperymentem rozstrzygającym tezę jest porównanie dwóch
\textbf{reguł decyzyjnych} na tych samych modelach. Reguła statyczna alarmuje,
gdy PD najnowszego okna przekracza próg ($\mathrm{PD}_{W3} \geq \theta$);
reguła monitorująca --- gdy próg przekracza PD któregokolwiek okna
($\max(\mathrm{PD}_{W0..W3}) \geq \theta$). Ponieważ druga reguła jest
nadzbiorem pierwszej (przy tym samym progu flaguje co najmniej tyle samo osób,
ale też generuje więcej fałszywych alarmów), uczciwe porównanie wymaga
wyrównania budżetu fałszywych alarmów: progi $\theta$ dobierane są niezależnie
dla każdej reguły tak, aby odsetek fałszywych alarmów (FA) wśród spłacających
wynosił zadane 5, 10 lub 20\%.

\begin{table}[htbp]
\centering
\caption{Zmiana czułości (monitoring $-$ statyka, pp) przy zadanych budżetach FA
oraz wartości specyficzne dla monitoringu (przy FA = 10\%).}
\label{tab:staticdyn}
\begin{tabular}{lrrrrr}
\toprule
Model & $\Delta$ FA=5\% & $\Delta$ FA=10\% & $\Delta$ FA=20\% & Tylko-monitoring & Śr. wyprzedzenie \\
\midrule
Random Forest & $+6{,}9$ & $\mathbf{-10{,}9}$ & $+2{,}3$ & 48 & 1,96 okna \\
XGBoost       & $+8{,}3$ & $\mathbf{-5{,}1}$  & $+0{,}6$ & 47 & 2,05 \\
LightGBM      & $+0{,}7$ & $\mathbf{-5{,}0}$  & $-2{,}0$ & 52 & 2,06 \\
CatBoost      & $+8{,}1$ & $\mathbf{-6{,}4}$  & $-8{,}6$ & 39 & 2,09 \\
\textbf{LSTM} & $+10{,}0$ & $\mathbf{+2{,}6}$ & $-2{,}3$ & \textbf{74} & 2,04 \\
\bottomrule
\end{tabular}
\end{table}

Wyniki są mieszane i~praca raportuje je bez upiększeń. Przy kanonicznym
budżecie FA~=~10\% reguła statyczna osiąga wyższą czułość dla wszystkich
czterech modeli statycznych (od $-5{,}0$ do $-10{,}9$~pp): agregator maksimum
po czterech silnie skorelowanych oknach ,,widzi'' czterokrotnie więcej szumu
niż pojedyncza, skalibrowana ocena W3, więc wyrównanie budżetu fałszywych
alarmów wymusza wyżej położony próg i~utratę części wykryć. Wyjątkiem jest
\textbf{LSTM --- jedyny model sekwencyjny --- dla którego reguła monitorująca
wygrywa także na czułości ($+2{,}6$~pp)}, co jest spójne z~hipotezą, że
architektura przetwarzająca sekwencję najlepiej wykorzystuje informację zawartą
w~trajektorii.

Wartość monitoringu ujawnia się w~dwóch wymiarach nieuchwytnych dla samej
czułości. Po pierwsze, \textbf{wyprzedzenie czasowe}: wśród wykrytych
defaultujących alarm pada średnio około dwa okna przed W3 --- w~praktyce
bankowej to czas na działania miękkie (kontakt, restrukturyzacja, obniżenie
limitu), których wartość nie jest ujęta w~macierzy pomyłek. Po drugie,
\textbf{wykrycia unikalne}: 39--74 defaultujących na model, których reguła
statyczna nie flaguje w~ogóle --- klienci, u~których ryzyko uwidoczniło się we
wcześniejszych oknach i~,,uspokoiło'' w~W3. Reguła monitorująca jest zatem
komplementarna wobec statycznej, nie substytucyjna; rozstrzygnięcie bilansu
zależy od modelu kosztów instytucji (im droższe późne wykrycie względem
fałszywego alarmu, tym większa wartość monitoringu).

Interpretację uzupełnia diagnoza rozkładu PD per okno, przeprowadzona osobno
dla obu klas. W~modelach drzewiastych średnie PD klientów spłacających jest
praktycznie identyczne we wszystkich oknach (różnica W0~$-$~W3 od $-0{,}002$ do
$-0{,}004$, częstość alarmów stabilna), a~u~defaultujących rośnie wyraźnie ku
W3 (o~0,065--0,071). Oznacza to, że (a)~modele te nie zawyżają ocen na
starszych oknach --- \textbf{nie występuje istotne przesunięcie rozkładu}
między treningiem na W3 a~inferencją na W0..W2, oraz (b)~dominacja najstarszego
okna w~histogramie pierwszych alertów jest artefaktem zliczania pierwszego
przekroczenia przy zbliżonej częstości alarmów per okno, nie dowodem na masowe
,,wczesne sygnały''; sygnał narastania ryzyka u~defaultujących jest natomiast
rzeczywisty. Jedynie LSTM wykazuje niewielki dryf również dla klasy spłacającej
($+0{,}004$ średniego PD i~45,4\% wobec 39,6\% częstości alarmów na W0 względem
W3) --- jego wynik w~regule monitorującej należy więc czytać z~tą poprawką,
choć porównanie przy stałym budżecie fałszywych alarmów
z~tabeli~\ref{tab:staticdyn} ten efekt z~konstrukcji neutralizuje.

\begin{figure}[htbp]
\centering
\includegraphics[width=.8\textwidth]{static_vs_dynamic_catboost_w3}
\caption{Reguła statyczna vs monitorująca --- krzywe czułość/fałszywe alarmy
(CatBoost; komplet pięciu modeli w~repozytorium). Źródło: opracowanie własne.}
\label{fig:staticdyn}
\end{figure}

\section{Interpretowalność modeli i ich przydatność decyzyjna}

Wyjaśnialność zrealizowano metodą SHAP (rozdz.~2.5): dla czterech modeli
drzewiastych system zwraca przy każdej predykcji pięć cech o~największym
wkładzie, z~konwencją znaku ,,wartość dodatnia podnosi PD''. Wyjaśnienia
liczone są przez TreeExplainer na estymatorze bazowym (sprzed kalibracji) ---
korekta izotoniczna jest monotoniczna, więc nie zmienia rankingu cech,
natomiast wartości SHAP opisują skalę modelu bazowego i~nie sumują się do
serwowanego, skalibrowanego PD; jest to udokumentowany kompromis na rzecz czasu
odpowiedzi (pojedyncze wyjaśnienie $\sim$100~ms, wobec limitu 2~s). Dla LSTM
wyjaśnienia pominięto (TreeExplainer nie ma zastosowania, a~KernelExplainer
przekracza budżet czasowy). Rankingi cech są zbieżne między modelami
i~z~ważnościami lasu losowego: dominują świeże zachowania płatnicze
(\texttt{PAY\_max}, \texttt{recent\_pay\_status}, \texttt{late\_count},
\texttt{utilization\_rate}), zmienne demograficzne zajmują odległe pozycje ---
struktura zależności w~danych, nie artefakt jednej rodziny algorytmów.

\subsection{Audyt sprawiedliwości względem atrybutu chronionego}

Audyt przeprowadzono biblioteką fairlearn~\cite{bird2020} dla pięciu
skalibrowanych modeli, względem atrybutu SEX, na zbiorze testowym (6~000
obserwacji; 2~402 mężczyzn, 3~598 kobiet), przy binaryzacji progami kosztowymi
z~sekcji~4.7 --- audyt odzwierciedla więc faktyczny punkt pracy systemu, a~nie
umowny próg 0,5. Zastosowano dwie komplementarne miary. Różnica parytetu
demograficznego (DPD) porównuje częstość alarmów między grupami, abstrahując od
rzeczywistych zdarzeń; różnica wyrównanych szans (EOD)~\cite{hardt2016}
warunkuje na prawdziwej etykiecie i~mierzy większą z~dwóch luk: wykrywalności
faktycznych defaultów (TPR) oraz fałszywych alarmów wśród spłacających (FPR).

Rozróżnienie to ma znaczenie zasadnicze. Częstość defaultu różni się między
grupami (23,4\% mężczyźni wobec 21,3\% kobiety w~zbiorze testowym), wskutek
czego nawet bezbłędny klasyfikator wykazywałby DPD około $+0{,}021$; wymuszenie
pełnego parytetu demograficznego wymagałoby traktowania osób o~identycznym
profilu ryzyka odmiennie ze względu na płeć --- co samo stanowiłoby
dyskryminację bezpośrednią i~stałoby w~sprzeczności z~obowiązkiem rzetelnej
oceny zdolności kredytowej (art.~70 Prawa bankowego). Dlatego DPD raportowane
jest jako miara przesiewowa, a~kryterium oceny stanowią wyrównane szanse ---
równy rozkład ciężaru błędów przy danym stanie faktycznym.

\begin{table}[htbp]
\centering
\caption{Audyt fairness przy progach kosztowych (limit DoD: 0,10).}
\label{tab:fairness}
\begin{tabular}{lrrrl}
\toprule
Model & Próg & DPD & EOD & Składnik dominujący EOD \\
\midrule
Random Forest & 0,145 & $+0{,}035$ & $+0{,}028$ & luka FPR (0,375 vs 0,347) \\
XGBoost       & 0,165 & $+0{,}036$ & $+0{,}028$ & luka FPR (0,369 vs 0,341) \\
LightGBM      & 0,160 & $+0{,}035$ & $+0{,}027$ & luka FPR (0,371 vs 0,344) \\
CatBoost      & 0,160 & $+0{,}039$ & $+0{,}033$ & luka FPR (0,394 vs 0,361) \\
LSTM          & 0,155 & $\mathbf{+0{,}006}$ & $+0{,}021$ & luka TPR, odwrócona (0,727 vs 0,748) \\
\bottomrule
\end{tabular}
\end{table}

\begin{figure}[htbp]
\centering
\includegraphics[width=.95\textwidth]{fig_5_10_audyt_fairness}
\caption{Audyt fairness: DPD/EOD z~limitem i~luką strukturalną, selection rate
oraz TPR/FPR per grupa. Źródło: opracowanie własne.}
\label{fig:fairness}
\end{figure}

Wszystkie modele spełniają kryterium $|\mathrm{DPD}|, |\mathrm{EOD}| \leq
0{,}10$ z~co najmniej trzykrotnym marginesem --- hipoteza H3 jest potwierdzona.
Interpretując wyniki na tle luki strukturalnej $+0{,}021$: modele drzewiaste
dokładają od siebie 0,014--0,018 ponad różnicę base rate, a~lukę wyrównanych
szans tworzy w~nich konsekwentnie składnik fałszywych alarmów --- rzetelnie
spłacający mężczyźni są błędnie flagowani o~2,7--3,3~pp częściej niż kobiety.
Na szczególną uwagę zasługuje LSTM: jako jedyny nie otrzymuje na wejściu
żadnych cech demograficznych (tensor $3\times3$ zawiera wyłącznie sekwencję
statusów, sald i~wpłat) i~jako jedyny osiąga DPD poniżej luki strukturalnej,
z~odwróconym znakiem luk błędów.

Kontrast ten podsunął naturalny eksperyment ablacyjny, który przeprowadzono
wprost: cztery modele drzewiaste przetrenowano w~identycznym protokole
\textbf{z~usuniętymi kolumnami SEX} (atrybut służy wtedy wyłącznie do pomiaru
metryk). Rezultat: zdolność predykcyjna nie zmieniła się w~granicach szumu
($|\Delta\mathrm{AUC}| \leq 0{,}001$ dla każdego modelu), natomiast parytet
uległ wyraźnej poprawie --- DPD XGBoost spadło z~$+0{,}036$ do $+0{,}011$,
LightGBM z~$+0{,}035$ do $+0{,}007$, lasu losowego z~$+0{,}035$ do $+0{,}027$,
CatBoost z~$+0{,}039$ do $+0{,}022$, z~analogicznym spadkiem EOD (XGBoost do
$+0{,}004$, LightGBM do $+0{,}002$). Wniosek jest dwojaki. Po pierwsze,
umiarkowane luki modeli pełnocechowych wynikały w~istotnej części
z~bezpośredniego użycia atrybutu chronionego, nie wyłącznie z~korelatów ---
a~jego usunięcie \textbf{nic nie kosztuje}: to jednoznaczna rekomendacja
wdrożeniowa (w~systemie produkcyjnym zmienna płci powinna zostać usunięta
z~wektora cech, pozostając w~danych wyłącznie do celów audytu, na co wprost
pozwala art.~10 ust.~5 rozporządzenia o~sztucznej inteligencji [25]).
Po drugie, wariant badawczy z~jawnym SEX spełnił swoją rolę: pozwolił
skwantyfikować wpływ atrybutu zamiast zakładać go a~priori.

Wobec spełnienia kryterium H3 z~zapasem, aktywna mitygacja nie jest konieczna;
zasadne jest podejście proporcjonalne: audyt DPD/EOD jako stała bramka
wydaniowa po każdym retreningu, usunięcie SEX z~cech jako krok wdrożeniowy,
a~jako udokumentowany mechanizm warunkowy --- korekta post-hoc progów metodą
ThresholdOptimizer (fairlearn) z~ograniczeniem wyrównanych szans, która nie
narusza skalibrowanych prawdopodobieństw (rdzenia trajektorii PD) i~wpina się
w~istniejącą, progową warstwę decyzyjną. Jej zastosowanie w~systemie wczesnego
ostrzegania --- gdzie alert kieruje do przeglądu analityka, nie do decyzji
odmownej --- wyrównuje ciężar bycia objętym przeglądem, a~nie warunki oferty.

\section{Dyskusja wyników i weryfikacja hipotez badawczych}

\textbf{H1 --- potwierdzona.} Modele 3-miesięczne zachowują jakość
6-miesięcznych baseline'ów: las losowy 0,7741 wobec 0,7792 ($-0{,}51$~pp),
XGBoost 0,7761 wobec 0,7818 ($-0{,}57$~pp), LSTM 0,7614 wobec 0,7686
($-0{,}72$~pp) --- wszystkie straty poniżej 1~pp (porównanie orientacyjne:
baseline'y nie były kalibrowane). Skrócenie okna, warunkujące istnienie
trajektorii, nie odbiera modelom zdolności dyskryminacyjnej.

\textbf{H2 --- potwierdzona częściowo, w~brzmieniu przyjętym w~3.2.}
Wcześniejsza detekcja: tak --- średnio $\sim$2~okna wyprzedzenia; wykrycia
niedostępne regule statycznej: tak --- 39--74 na model. Natomiast czułość przy
stałym budżecie fałszywych alarmów 10\% jest dla czterech modeli statycznych
niższa (do $-10{,}9$~pp) --- z~wyjątkiem LSTM ($+2{,}6$~pp). Teza obroniona
w~brzmieniu ,,wcześniej i~komplementarnie'', nie ,,czulej''; dodatkowo wynik
LSTM wskazuje, że pełne wykorzystanie trajektorii wymaga architektury
sekwencyjnej.

\textbf{H3 --- potwierdzona.} Maksymalna wartość $|\mathrm{DPD}|/|\mathrm{EOD}|$
w~całej piątce wynosi 0,039 przy limicie 0,10; kontr-eksperyment bez SEX domyka
analizę rekomendacją wdrożeniową.

\textbf{Ograniczenia.} (1)~Jeden zbiór (Tajwan 2005) i~jeden podział ---
różnice między modelami drzewiastymi raportowano z~przedziałami bootstrap;
generalizacja na inne portfele niezweryfikowana. (2)~Monitoring symulowany
retrospektywnie na danych przekrojowych: okna nakładają się, a~horyzont
predykcji różni się między oknami (rozdz.~3.3.4). (3)~Progi kosztowe wyznaczane
na części kalibracyjnej wspólnej z~kalibratorem (kompromis opisany w~4.7).
(4)~Wyjaśnienia SHAP operują na skali modeli bazowych, nie skalibrowanej.
(5)~Model kosztu $\mathrm{FN} = 5 \times \mathrm{FP}$ jest założeniem
eksperckim --- wnioski sekcji~5.4 są funkcją tej asymetrii.
